%% Hierarchical timing wheels: the general algorithm, with Tokio 1.54.1 as the reference implementation.

\section{Hierarchical timing wheels}
\label{sec:timing-wheels}

\textit{A runtime, a kernel or a network stack may hold millions of timers at once, almost all of which will be cancelled before they fire. A timing wheel makes every timer operation take constant time, however many timers exist.}

A timer service needs four operations: \emph{insert} a timer, \emph{cancel} one, find the \emph{next} deadline so its owner knows how long to sleep, and \emph{expire} the timers that are due. The obvious structures each make one of them slow:

\begin{itemize}
    \item \textbf{An unsorted list} inserts in constant time, but finding the next deadline means scanning every timer.
    \item \textbf{A sorted list} finds the next deadline at its head, but inserting means walking to the right place.
    \item \textbf{A binary heap} does both in $O(\log n)$, and cancelling needs each timer to know its position in the heap.
\end{itemize}

A timing wheel trades exactness of order for speed. Time is cut into fixed \emph{ticks}, and timers are hashed by deadline into slots, like the hours on a clock face. A \emph{hierarchical} wheel stacks several such clocks, each coarser than the one below, so a few small arrays cover deadlines years away. The design is due to Varghese and Lauck.\footnote{G.~Varghese and A.~Lauck, ``Hashed and Hierarchical Timing Wheels'', \emph{IEEE/ACM Transactions on Networking}, 1997.}

This section uses the wheel in the Tokio runtime, version 1.54.1, as its reference implementation: every algorithm is read from \texttt{runtime/time/wheel/} and \texttt{runtime/time/entry.rs}, and a grey ``Tokio:'' line names the functions it comes from. How Tokio's runtime drives the wheel is in Section~\ref{sec:wakers-ch15}.

\subsection{The structure}

\begin{definition}{Hierarchical timer wheel}{timer-wheel}
    \index{Timer wheel}
    A \emph{hierarchical timer wheel}, as Tokio builds it, has six levels of 64 slots. A level-0 slot spans 1 millisecond, and each level above is 64 times coarser: 64\,ms, about 4.1\,s, 4.4\,min, 4.7\,h and 12.4 days. Together the levels span $64^6$ milliseconds, a little over two years. Each level also keeps a 64-bit mask recording which of its slots are occupied.
\end{definition}

\subsection{Ticks and the wheel's clock}

\begin{itemize}
    \item The wheel counts time in integer \emph{ticks}. In Tokio a tick is one millisecond.
    \item A deadline is rounded \emph{up} to a whole tick, so a timer never fires early.
    \item The wheel keeps its own clock, \code{elapsed}: the tick up to which it has processed timers. It only moves forward, and only inside the advance algorithm below. Every placement decision is made relative to \code{elapsed}, not to the wall-clock ``now''.
\end{itemize}

At level $\ell$, a slot spans $64^{\ell}$ ticks and the whole level spans $64^{\ell+1}$. Writing a tick as a 36-bit number split into six 6-bit groups, group $\ell$ is exactly the slot index at level $\ell$.

\subsection{What the wheel stores}

\begin{description}
    \item[Wheel] \code{elapsed}; six \code{Level}s; and a \code{pending} list of timers that are due but not yet fired.
    \item[Level] An array of 64 slots, each an \emph{intrusive} doubly linked list (the links live inside the timer entries themselves, so adding or unlinking an entry allocates nothing), plus \code{occupied}, a 64-bit mask whose bit $s$ is set exactly when slot $s$ is non-empty.
    \item[Timer entry] Owned by whoever is waiting (in Tokio, a \code{Sleep}) and linked into the wheel. Besides what to do when it fires (in Tokio, a waker to call), it carries two deadlines, defined next.
\end{description}

\begin{definition}{Filing deadline, true deadline}{filing-true-deadline}
    \index{Timer wheel!filing deadline}
    A timer's \emph{true deadline} (in Tokio, the value in the entry's atomic \code{state}) is when it should fire. Its \emph{filing deadline} (\code{registered\_when}) is the deadline the wheel used to choose its level and slot, or the marker \textsc{Deregistered} once it has moved to the pending list. The two are equal when a timer is inserted; they differ only after the deadline is moved later (Lemma~\ref{lem:lazy-extension}). The pseudocode below writes them as $x.\mathit{when}$ and $x.\mathit{filed}$.
\end{definition}

\subsection{Insertion: choosing a level and a slot}

The level is the 6-bit group of the highest bit in which the wheel's clock and the deadline differ; the slot is the deadline's value in that group. Figure~\ref{fig:tw-level-for} works one case, and Algorithm~\ref{alg:tw-level-for} gives the rule.

\begin{figure}[htbp]
    \centering
    \figurefonts\resizebox{\linewidth}{!}{\begin{tikzpicture}[font=\sffamily\footnotesize, >={Stealth},
        grp/.style={draw=gray!50, fill=white, minimum width=1.5cm, minimum height=0.6cm, font=\ttfamily\footnotesize},
        same/.style={grp, fill=gray!12, text=gray},
        diff/.style={grp, draw=StmtProposition, fill=orange!12, line width=0.8pt},
        rowlbl/.style={anchor=east, font=\sffamily\footnotesize}]
        \foreach \g [count=\i from 0] in {5,4,3,2,1,0}{
            \node[text=gray, font=\sffamily\scriptsize] at ({1.7*\i},0.65) {group \g};
        }
        % elapsed = 4096
        \node[rowlbl] at (-1.0,0) {\code{elapsed} $= 4096$};
        \foreach \v [count=\i from 0] in {0,0,0,1}{\node[same] at ({1.7*\i},0) {\v};}
        \node[diff] at (6.8,0) {0};
        \node[grp] at (8.5,0) {0};
        % when = 5000
        \node[rowlbl] at (-1.0,-0.9) {deadline $= 5000$};
        \foreach \v [count=\i from 0] in {0,0,0,1}{\node[same] at ({1.7*\i},-0.9) {\v};}
        \node[diff] (slot) at (6.8,-0.9) {14};
        \node[grp] at (8.5,-0.9) {8};
        % xor | 63
        \node[rowlbl] at (-1.0,-1.8) {XOR, then $\mathbin{|}\,63$};
        \foreach \v [count=\i from 0] in {0,0,0,0}{\node[same] at ({1.7*\i},-1.8) {\v};}
        \node[diff] at (6.8,-1.8) {14};
        \node[grp] at (8.5,-1.8) {63};
        % annotations
        \draw[gray] (-0.75,-2.35) -- (5.85,-2.35);
        \node[text=gray, align=center] at (2.55,-2.7) {groups 5--2 agree: same level-1 window};
        \draw[->, StmtProposition, thick] (6.8,-2.15) -- (6.8,-2.9) node[below, align=center, text=StmtProposition, font=\sffamily\footnotesize\bfseries] {highest differing group is 1:\\ level 1, slot 14};
        \node[anchor=west, text=gray, text width=2.3cm] at (9.35,-1.8) {forced on: the level is never below 0};
    \end{tikzpicture}}
    \caption{Choosing a level and slot for a deadline of 5,000 ticks when the wheel's clock is at 4,096, as in step 6 of Table~\ref{tab:tw-example}. Each box is one 6-bit group of the tick, in decimal.}
    \label{fig:tw-level-for}
\end{figure}

\begin{algorithm}[htbp]
    \caption{Choosing a level and a slot.}
    \label{alg:tw-level-for}
    \begin{algorithmic}[1]
        \Statex {\color{gray!90!black}\footnotesize Tokio: \texttt{level\_for}, \texttt{slot\_for}}
        \Function{LevelFor}{$e, w$}\Comment{$e$: wheel clock; $w$: deadline; both in ticks}
            \State $m \gets (e \oplus w) \mathbin{|} 63$\Comment{set bits 0--5, so the result is at least level 0}
            \If{$m \ge 2^{36}$}
                \State \Return $5$\Comment{more than $64^6$ ticks away: the top level}
            \EndIf
            \State \Return $\lfloor \log_2 m \rfloor \operatorname{div} 6$\Comment{group of the highest differing bit}
        \EndFunction
        \Statex
        \Function{SlotFor}{$w, \ell$}
            \State \Return $\lfloor w / 64^{\ell} \rfloor \bmod 64$\Comment{group $\ell$ of $w$}
        \EndFunction
    \end{algorithmic}
\end{algorithm}

\begin{algorithm}[htbp]
    \caption{Inserting a timer, under the driver lock.}
    \label{alg:tw-insert}
    \begin{algorithmic}[1]
        \Statex {\color{gray!90!black}\footnotesize Tokio: \texttt{Wheel::insert}, called from \texttt{reregister}}
        \Function{Insert}{$x$}
            \State $w \gets x.\mathit{when}$;\enspace $x.\mathit{filed} \gets w$
            \If{$w \le \mathit{elapsed}$}
                \State \Return \textsc{Elapsed}\Comment{already due: the caller fires it at once}
            \EndIf
            \State $\ell \gets \Call{LevelFor}{\mathit{elapsed}, w}$;\enspace $s \gets \Call{SlotFor}{w, \ell}$
            \State push $x$ onto the front of $\mathit{slot}_\ell[s]$
            \State $\mathit{occupied}_\ell \gets \mathit{occupied}_\ell \mathbin{|} 2^{s}$
            \If{$w < \mathit{next\_wake}$}\Comment{sooner than the sleeper expects}
                \State wake the thread sleeping until $\mathit{next\_wake}$
            \EndIf
        \EndFunction
    \end{algorithmic}
\end{algorithm}

The last step is for whoever drives the wheel. It sleeps until the next expiration, recorded as $\mathit{next\_wake}$, and a new timer that is due sooner must cut that sleep short. Section~\ref{sec:wakers-ch15} shows how Tokio's runtime does it.

\begin{lemma}{Placement}{wheel-placement}
    A timer filed at level $\ell$ below the top agrees with the wheel's clock in every group above $\ell$, and is greater in group $\ell$. It therefore lies in the current window of level $\ell$, in a slot strictly after the one that holds the clock.
    \why By Algorithm~\ref{alg:tw-level-for}, group $\ell$ holds the highest bit in which they differ, so all higher groups agree. The deadline is later than the clock, so at that highest differing bit the deadline has the 1; its group $\ell$, which is its slot, is therefore the larger. At level 0 the same holds, since the deadline differs from the clock only in group 0.
\end{lemma}

\begin{lemma}{Constant-time insert}{constant-time-insert}
    Adding a timer takes constant time, however many timers exist.
    \why Algorithm~\ref{alg:tw-insert} is an XOR, a leading-zero count, a shift, a push onto the front of a linked list and a bit set. None of it looks at other timers.
\end{lemma}

\subsection{Deletion}

Most timers never fire. A timeout guarding a request that succeeds in 3\,ms is cancelled long before its deadline, so cancellation (Algorithm~\ref{alg:tw-remove}) is at least as common as expiry, and must be just as cheap as insertion.

\begin{algorithm}[htbp]
    \caption{Removing a timer.}
    \label{alg:tw-remove}
    \begin{algorithmic}[1]
        \Statex {\color{gray!90!black}\footnotesize Tokio: \texttt{Wheel::remove}, from \texttt{clear\_entry} (a dropped \texttt{Sleep}) or \texttt{reregister}}
        \Function{Remove}{$x$}
            \If{$x.\mathit{filed} = \textsc{Deregistered}$}
                \State unlink $x$ from $\mathit{pending}$;\enspace \Return\Comment{due, but not yet fired}
            \EndIf
            \State $\ell \gets \Call{LevelFor}{\mathit{elapsed}, x.\mathit{filed}}$;\enspace $s \gets \Call{SlotFor}{x.\mathit{filed}, \ell}$
            \State unlink $x$ from $\mathit{slot}_\ell[s]$\Comment{doubly linked: no search}
            \If{$\mathit{slot}_\ell[s]$ is now empty}
                \State $\mathit{occupied}_\ell \gets \mathit{occupied}_\ell \mathbin{\&} \lnot 2^{s}$
            \EndIf
        \EndFunction
    \end{algorithmic}
\end{algorithm}

Removal stores nothing about where the timer was filed: it recomputes the level from the filing deadline and the wheel's clock \emph{now}. That is only correct because of the next lemma.

\begin{lemma}{Stable filing}{stable-filing}
    From the moment a timer is filed until its slot is processed, \textsc{LevelFor}$(\mathit{elapsed}, x.\mathit{filed})$ keeps returning the level it is filed on.
    \why The clock moves only in the advance algorithm (Algorithm~\ref{alg:tw-advance}): either to the start of the next due slot, or to ``now'' when nothing is due by then. Neither move passes the start of any occupied slot without processing it. So until the timer's slot is processed, the clock stays in the same level-$\ell$ window and before the timer's slot, and Lemma~\ref{lem:wheel-placement} goes on giving the same level.
\end{lemma}

\begin{lemma}{Constant-time removal}{constant-time-remove}
    Removing a timer takes constant time, however many timers exist.
    \why By Lemma~\ref{lem:stable-filing}, Algorithm~\ref{alg:tw-remove} finds the right slot without searching. Unlinking from a doubly linked list and clearing a bit take constant time.
\end{lemma}

\subsection{Finding the next deadline}

The driver needs one number before it sleeps: the tick at which the wheel next has work to do. Algorithm~\ref{alg:tw-next} finds it without looking at a single timer.

\begin{algorithm}[htbp]
    \caption{Finding the next expiration.}
    \label{alg:tw-next}
    \begin{algorithmic}[1]
        \Statex {\color{gray!90!black}\footnotesize Tokio: \texttt{Level::next\_occupied\_slot}, \texttt{Wheel::next\_expiration}}
        \Function{NextOccupiedSlot}{$\ell, e$}
            \If{$\mathit{occupied}_\ell = 0$}
                \State \Return \textsc{None}
            \EndIf
            \State $n \gets (\lfloor e / 64^{\ell} \rfloor \bmod 64) + 1$\Comment{start just after the clock's slot}
            \State $r \gets \operatorname{rotr}(\mathit{occupied}_\ell, n)$\Comment{rotate the 64-bit mask right by $n$}
            \State \Return $(\operatorname{tzcnt}(r) + n) \bmod 64$\Comment{first occupied slot, wrapping round}
        \EndFunction
        \Statex
        \Function{NextExpiration}{}
            \If{$\mathit{pending}$ is not empty}
                \State \Return $(0, 0, \mathit{elapsed})$\Comment{something is already due}
            \EndIf
            \For{$\ell \gets 0$ \textbf{to} $5$}\Comment{lowest level first}
                \State $s \gets \Call{NextOccupiedSlot}{\ell, \mathit{elapsed}}$
                \If{$s \ne \textsc{None}$}
                    \State $d \gets \mathit{elapsed}$ rounded down to a multiple of $64^{\ell+1}$, plus $s \cdot 64^{\ell}$
                    \If{$d \le \mathit{elapsed}$}
                        \State $d \gets d + 64^{\ell+1}$\Comment{top level only: the next rotation}
                    \EndIf
                    \State \Return $(\ell, s, d)$\Comment{$d$ is the slot's \emph{start}}
                \EndIf
            \EndFor
            \State \Return \textsc{None}
        \EndFunction
    \end{algorithmic}
\end{algorithm}

\begin{figure}[htbp]
    \centering
    \newcommand{\wkmask}[3]{% y, clock slot, caption text
        \foreach \s in {0,...,63}{\fill[rounded corners=0.3pt, gray!15] ({0.2*\s},#1) rectangle ++(0.16,0.45);}
        \foreach \s in {10,40}{\fill[rounded corners=0.3pt, StmtLemma] ({0.2*\s},#1) rectangle ++(0.16,0.45);}
        \draw[StmtProposition, line width=0.8pt, rounded corners=0.5pt] ({0.2*#2-0.03},{#1-0.03}) rectangle ++(0.22,0.51);
        \node[anchor=east, align=right] at (-0.2,{#1+0.22}) {#3};
    }
    \figurefonts\resizebox{\linewidth}{!}{\begin{tikzpicture}[font=\sffamily\footnotesize, >={Stealth}]
        \foreach \t in {0,10,20,30,40,50,63}{\node[text=gray, font=\sffamily\scriptsize] at ({0.2*\t+0.08},0.75) {\t};}
        \wkmask{0}{20}{clock in slot 20}
        \draw[->, StmtProposition, thick] ({0.2*21+0.08},-0.15) -- ({0.2*40+0.08},-0.15);
        \node[anchor=west, text=StmtProposition] at ({0.2*40+0.4},-0.15) {19 trailing zeros: slot $21 + 19 = 40$};
        \wkmask{-1.3}{41}{clock in slot 41}
        \draw[->, StmtProposition, thick] ({0.2*42+0.08},-1.45) -- (12.9,-1.45);
        \draw[->, StmtProposition, thick] (-0.05,-1.6) -- ({0.2*10+0.08},-1.6);
        \node[anchor=west, text=StmtProposition] at ({0.2*10+0.4},-1.65) {wraps round: slot $(42 + 32) \bmod 64 = 10$};
    \end{tikzpicture}}
    \caption{\textsc{NextOccupiedSlot} on a level whose slots 10 and 40 are occupied (green). Rotating the mask so that the slot after the clock's (orange) becomes bit 0, the count of trailing zeros is the distance to the next occupied slot. These are cases from Tokio's own tests.}
    \label{fig:tw-next-slot}
\end{figure}

Two details in Algorithm~\ref{alg:tw-next} deserve a word.

\begin{itemize}
    \item \textbf{The deadline returned is a slot's start, not a timer's.} At level 0 the two coincide. At a higher level, reaching the slot's start means it is time to cascade the slot's timers to finer levels, not to fire them.
    \item \textbf{The scan skips the clock's own slot.} By Lemma~\ref{lem:wheel-placement}, below the top level that slot is always empty. Only the top level, which acts as a ring for deadlines more than $64^6$ ticks away, can have a timer there, and that timer is a whole rotation in the future. (Tokio 1.54 made this explicit with the ``$+1$''.)
\end{itemize}

\begin{lemma}{Lowest level first}{lowest-level-first}
    If any timer is filed at level $\ell$, the next occupied slot of level $\ell$ starts before every occupied slot of every higher level. So the first level with an occupied slot gives the wheel's next expiration.
    \why By Lemma~\ref{lem:wheel-placement}, every level-$\ell$ timer lies inside the current level-$\ell$ window, which ends at the next multiple of $64^{\ell+1}$ after the clock. A timer at a higher level $\ell'$ is greater than the clock in group $\ell'$, so its slot starts at or after that same boundary. Tokio checks this property with an assertion in debug builds.
\end{lemma}

\begin{lemma}{Constant-time next deadline}{constant-time-next-deadline}
    Finding the next deadline takes at most one rotate and one trailing-zero count per level: six of each in all, however many timers exist.
    \why Algorithm~\ref{alg:tw-next} reads only the six occupancy masks, and by Lemma~\ref{lem:lowest-level-first} stops at the first non-empty one.
\end{lemma}

\subsection{Advancing time: firing and cascading}

Whenever whoever drives the wheel wakes up, it reads the clock and advances the wheel to it (Algorithm~\ref{alg:tw-advance}). Each due slot is taken off whole. Its timers that are due move to the pending list; the rest are filed again, relative to the slot's start.

\begin{algorithm}[htbp]
    \caption{Advancing the wheel to tick $t$.}
    \label{alg:tw-advance}
    \begin{algorithmic}[1]
        \Statex {\color{gray!90!black}\footnotesize Tokio: \texttt{Wheel::poll}, \texttt{Wheel::process\_expiration}}
        \Function{Advance}{$t$}
            \Loop
                \State $E \gets \Call{NextExpiration}{}$
                \If{$E = \textsc{None}$ \textbf{or} $E.d > t$}
                    \State $\mathit{elapsed} \gets t$;\enspace \textbf{break}\Comment{nothing more is due}
                \EndIf
                \State $(\ell, s, d) \gets E$
                \State $L \gets \mathit{slot}_\ell[s]$;\enspace $\mathit{slot}_\ell[s] \gets \emptyset$;\enspace clear bit $s$ of $\mathit{occupied}_\ell$
                \ForAll{$x \in L$}
                    \If{$x.\mathit{when} \le d$}
                        \State $x.\mathit{filed} \gets \textsc{Deregistered}$;\enspace move $x$ to $\mathit{pending}$\Comment{due}
                    \Else
                        \State $x.\mathit{filed} \gets x.\mathit{when}$;\enspace $\ell' \gets \Call{LevelFor}{d, x.\mathit{when}}$\Comment{cascade}
                        \State push $x$ onto $\mathit{slot}_{\ell'}[\Call{SlotFor}{x.\mathit{when}, \ell'}]$ and set its bit
                    \EndIf
                \EndFor
                \State $\mathit{elapsed} \gets d$
            \EndLoop
            \State \Return $\mathit{pending}$\Comment{the caller fires each one}
        \EndFunction
    \end{algorithmic}
\end{algorithm}

Taking the whole slot first matters at the top level: a timer more than a full rotation away is filed straight back into the same slot, and must not be processed again in the same pass.

\begin{corollary}{Bounded cascades}{bounded-cascades}
    A timer inserted at level $\ell$ is filed again at most $\ell$ times before it fires, unless its deadline is moved later or lies more than $64^6$ ticks ahead.
    \why When a level-$\ell$ slot starting at $d$ is processed, each of its timers lies in $[d, d + 64^{\ell})$, so it agrees with $d$ in every group from $\ell$ up: \textsc{LevelFor}$(d, \cdot)$ returns a level below $\ell$. Each cascade therefore goes down at least one level, and level 0 fires.
\end{corollary}

\subsection{A worked example}

At time zero, five timers are created, with deadlines of 3\,ms, 70\,ms, 300\,ms, 5\,s and 90\,s. Figure~\ref{fig:tw-example} shows where Tokio's placement rule puts them.

\begin{figure}[htbp]
    \centering
    % Coordinates follow a 774 x 308 grid; y grows downwards. Slot i of a level starts at x = 118 + 10i.
    \newcommand{\wklevel}[4]{% top y, name, slot size, occupied slots
        \node[anchor=north west, inner sep=0] at (8,{#1+6}) {\small\textbf{#2}};
        \node[anchor=north west, inner sep=0, text=gray] at (8,{#1+22}) {#3};
        \foreach \t in {0,16,32,48,63}{\node[text=gray, font=\sffamily\tiny] at ({122+10*\t},{#1-6}) {\t};}
        \foreach \s in {0,...,63}{\fill[rounded corners=0.5pt, gray!15] ({118+10*\s},#1) rectangle ++(8,26);}
        \foreach \s in {#4}{\fill[rounded corners=0.5pt, StmtLemma] ({118+10*\s},#1) rectangle ++(8,26);}
        \draw[StmtProposition, line width=0.9pt, rounded corners=1pt] (116,{#1-2}) rectangle ++(12,30);
    }
    \figurefonts\resizebox{\linewidth}{!}{\begin{tikzpicture}[x=0.02cm, y=-0.02cm, font=\sffamily\scriptsize]
        \wklevel{20}{Level 0}{64 $\times$ 1 ms}{3}
        \wklevel{116}{Level 1}{64 $\times$ 64 ms}{1,4}
        \wklevel{212}{Level 2}{64 $\times$ 4.1 s}{1,21}
        % timer labels
        \draw[StmtLemma] (152,47) -- (152,58);
        \node[anchor=north, align=center, inner sep=1pt] at (152,58) {\textcolor{StmtLemma}{\textbf{3 ms}}\\ \textcolor{gray}{\tiny slot 3}};
        \draw[StmtLemma] (132,143) -- (132,154);
        \node[anchor=north, align=center, inner sep=1pt] at (132,154) {\textcolor{StmtLemma}{\textbf{70 ms}}\\ \textcolor{gray}{\tiny slot 1}};
        \draw[StmtLemma] (162,143) -- (162,150) -- (180,150);
        \node[anchor=west, inner sep=1pt] at (180,150) {\textcolor{StmtLemma}{\textbf{300 ms}} \textcolor{gray}{\tiny $\cdot$ slot 4}};
        \draw[StmtLemma] (132,239) -- (132,250);
        \node[anchor=north, align=center, inner sep=1pt] at (132,250) {\textcolor{StmtLemma}{\textbf{5 s}}\\ \textcolor{gray}{\tiny slot 1}};
        \draw[StmtLemma] (332,239) -- (332,250);
        \node[anchor=north, align=center, inner sep=1pt] at (332,250) {\textcolor{StmtLemma}{\textbf{90 s}}\\ \textcolor{gray}{\tiny slot 21}};
    \end{tikzpicture}}
    \caption{Five timers placed at time zero. The first three of the six levels are shown; the orange outline marks the current slot on each level.}
    \label{fig:tw-example}
\end{figure}

Table~\ref{tab:tw-example} follows time forward. Each row is one call to \textsc{Advance} by a caller that sleeps until the next expiration in between; ``Timeout'' is how long it sleeps. Each row touches only the single slot that came due.

\begin{table}[htbp]
    \centering\small
    \begin{tabularx}{\linewidth}{r r Y}
        \toprule
        \textbf{Time} & \textbf{Timeout} & \textbf{What \textsc{Advance} does} \\
        \midrule
        3 ms & 3 ms & Level-0 slot 3 due: \textbf{fire} the 3 ms timer \\
        64 ms & 61 ms & Level-1 slot 1 due: cascade the 70 ms timer into level-0 slot 6 \\
        70 ms & 6 ms & \textbf{Fire} the 70 ms timer \\
        256 ms & 186 ms & Level-1 slot 4 due: cascade the 300 ms timer into level-0 slot 44 \\
        300 ms & 44 ms & \textbf{Fire} the 300 ms timer \\
        4,096 ms & 3,796 ms & Level-2 slot 1 due: cascade the 5 s timer into level-1 slot 14 \\
        4,992 ms & 896 ms & Level-1 slot 14 due: cascade it into level-0 slot 8 \\
        5,000 ms & 8 ms & \textbf{Fire} the 5 s timer \\
        86,016 ms & 81,016 ms & Level-2 slot 21 due: cascade the 90 s timer into level-1 slot 62 \\
        89,984 ms & 3,968 ms & Cascade it into level-0 slot 16 \\
        90,000 ms & 16 ms & \textbf{Fire} the 90 s timer \\
        \bottomrule
    \end{tabularx}
    \caption{The example run to completion, computed with Tokio's rule. A run of 10,030 random timers, checked the same way, fired every one exactly on its deadline.}
    \label{tab:tw-example}
\end{table}

\begin{itemize}
    \item Ninety seconds of timekeeping for five timers costs eleven brief wake-ups.
    \item Between them, the caller sleeps; in Tokio, that is the worker on driver duty, at no CPU cost (Axiom~\ref{ax:free-blocking}).
    \item Adding ten thousand more timers adds work only at the moments when some of them are due.
\end{itemize}

\subsection{Moving a deadline}

Timers are often moved. Servers do this constantly: an idle timeout is pushed back every time a packet arrives on the connection. Moving a deadline later is by far the common case, and a wheel can make it almost free (Algorithm~\ref{alg:tw-reset}).

\begin{algorithm}[htbp]
    \caption{Moving a timer's deadline to $w'$.}
    \label{alg:tw-reset}
    \begin{algorithmic}[1]
        \Statex {\color{gray!90!black}\footnotesize Tokio: \texttt{TimerEntry::reset}, \texttt{StateCell::extend\_expiration}}
        \Function{Reset}{$x, w'$}
            \If{$x$ has not fired \textbf{and} $w' \ge x.\mathit{when}$}
                \State $x.\mathit{when} \gets w'$\Comment{one compare-and-swap: no lock, wheel untouched}
                \State \Return
            \EndIf
            \State lock the wheel\Comment{an earlier deadline, or the timer already fired}
            \State \Call{Remove}{$x$};\enspace $x.\mathit{when} \gets w'$
            \If{\Call{Insert}{$x$} $=$ \textsc{Elapsed}}
                \State fire $x$ and call its waker
            \EndIf
        \EndFunction
    \end{algorithmic}
\end{algorithm}

\begin{lemma}{Lazy extension}{lazy-extension}
    \index{Timer wheel!lazy extension}
    Moving a timer's deadline later costs one atomic compare-and-swap, with no lock and no change to the wheel. The timer still fires on time; the price is at most one extra re-filing, paid when its old slot comes due.
    \why The timer stays filed under its old, earlier filing deadline, so its slot comes due no later than the new deadline: it can never fire late. When that slot is processed, Algorithm~\ref{alg:tw-advance} finds $x.\mathit{when} > d$ and files the timer again by its true deadline. From there it behaves like a timer inserted at that moment.
\end{lemma}

The model used to check Table~\ref{tab:tw-example} shows it on one timer. Insert a 70\,ms timer at tick 0: it is filed at level 1, slot 1. At tick 10, reset it to 300\,ms: nothing in the wheel changes. Then:

\begin{table}[htbp]
    \centering\small
    \begin{tabularx}{\linewidth}{r Y}
        \toprule
        \textbf{Time} & \textbf{What the driver does} \\
        \midrule
        64 ms & Level-1 slot 1 due. The true deadline (300) is later than the slot's start (64), so the timer is filed again, at level 1, slot 4: the one extra step \\
        256 ms & Level-1 slot 4 due: cascade into level-0 slot 44, exactly as for a timer created for 300\,ms \\
        300 ms & \textbf{Fire} \\
        \bottomrule
    \end{tabularx}
    \caption{A 70\,ms timer reset to 300\,ms at tick 10. The reset took no lock; the timer cost one extra re-filing at 64\,ms.}
    \label{tab:tw-lazy}
\end{table}

Moving a deadline \emph{earlier} cannot be lazy: the timer is filed too late to fire on time. The wheel's owner then takes the lock, removes the timer and inserts it again (Algorithms~\ref{alg:tw-remove} and \ref{alg:tw-insert}), both in constant time, and the insert may cut a sleeper short.

\begin{table}[htbp]
    \centering\small
    \begin{tabularx}{\linewidth}{>{\raggedright\arraybackslash}p{3.4cm} Y >{\raggedright\arraybackslash}p{3.0cm}}
        \toprule
        \textbf{Operation} & \textbf{Cost} & \textbf{Why} \\
        \midrule
        Insert & $O(1)$, lock & Lemma~\ref{lem:constant-time-insert} \\
        Remove (cancel) & $O(1)$, lock & Lemma~\ref{lem:constant-time-remove} \\
        Move the deadline later & $O(1)$, no lock & Lemma~\ref{lem:lazy-extension} \\
        Move the deadline earlier & $O(1)$, lock & Remove, then insert \\
        Next deadline & six rotates and trailing-zero counts at most & Lemma~\ref{lem:constant-time-next-deadline} \\
        Advance & $O(1)$ per due slot, plus $O(1)$ per timer in it & Algorithm~\ref{alg:tw-advance} \\
        Re-filings per timer & at most its starting level & Corollary~\ref{cor:bounded-cascades} \\
        \bottomrule
    \end{tabularx}
    \caption{The cost of every timer-wheel operation. None depends on how many timers exist.}
    \label{tab:tw-costs}
\end{table}

\begin{quotebox}{Summary}
    \begin{itemize}
        \item A hierarchical wheel files each timer at the 6-bit group of the highest bit in which its deadline differs from the wheel's clock (Algorithm~\ref{alg:tw-level-for}).
        \item Insert, remove and next-deadline are all constant time (Lemmas~\ref{lem:constant-time-insert}, \ref{lem:constant-time-remove}, \ref{lem:constant-time-next-deadline}). Removal can recompute a timer's slot because filings are stable (Lemma~\ref{lem:stable-filing}).
        \item Advancing processes one slot at a time: due timers fire, the rest cascade down, at most once per level (Corollary~\ref{cor:bounded-cascades}).
        \item Moving a deadline later is lock-free and lazy (Lemma~\ref{lem:lazy-extension}).
    \end{itemize}
\end{quotebox}

\begin{Exercise}[title={Timing wheels}]
    \begin{enumerate}
        \item At time zero, a 200\,ms timer is created. Using Algorithm~\ref{alg:tw-level-for}, on which level and slot of the timer wheel is it filed? When is it cascaded, into which slot, and when does it fire?
        \item A level-0 occupancy mask has slots 5 and 60 set, and the wheel's clock is at tick 130. Run \textsc{NextOccupiedSlot} and \textsc{NextExpiration} (Algorithm~\ref{alg:tw-next}) by hand: which slot do they find, and at what tick does it come due?
        \item Algorithm~\ref{alg:tw-advance} takes a due slot's whole list off the wheel before looking at any of its timers. What would go wrong if it processed the timers one by one, straight from the slot?
        \item A binary heap is another way to hold timers. Compare its costs with the wheel's for insert, cancel and finding the next deadline. When might a heap still be the better choice?
    \end{enumerate}
\end{Exercise}
\begin{Answer}
    \begin{enumerate}
        \item $200 = 3 \times 64 + 8$. ``Now'' (0) and the deadline first differ in the level-1 group of 6 bits, so the timer is filed on level 1, slot 3 (Algorithm~\ref{alg:tw-level-for}). When level-1 slot 3's window opens at $3 \times 64 = 192$\,ms, it is cascaded into level-0 slot $200 \bmod 64 = 8$, and it fires at 200\,ms. This matches the pattern of Table~\ref{tab:tw-example}.
        \item The scan starts at $n = (130 \bmod 64) + 1 = 3$. Rotating the mask right by 3 moves slot 5 to bit 2 and slot 60 to bit 57, so there are 2 trailing zeros and the slot is $(2 + 3) \bmod 64 = 5$. The clock rounded down to a multiple of 64 is 128, so the slot comes due at tick $128 + 5 = 133$. Slot 60 is not next, even though it is ``after'' slot 5 in the array: the scan starts from the clock's slot (2), not from slot 0.
        \item At the top level, a timer more than a full rotation ($64^6$ ticks) away is filed back into the very slot being processed. Processing the slot in place would meet that timer again in the same pass and file it back again, for ever. Taking the whole list first means each timer is looked at once per pass.
        \item A heap of $n$ timers inserts and cancels in $O(\log n)$ (cancelling needs each timer to record its index in the heap) and finds the next deadline in $O(1)$; the wheel does all three in constant time (Table~\ref{tab:tw-costs}). A heap can still be the better choice when there are few timers, when deadlines need finer resolution than a fixed tick, or when exact ordering of deadlines matters: it never rounds deadlines to ticks and never cascades.
    \end{enumerate}
\end{Answer}
